\documentclass[a4paper,11pt]{article}

% 重複していたパッケージの読み込みを1つに統合しました
\usepackage{amsmath, amssymb, amsthm}
\usepackage{tikz-cd}
\usepackage{booktabs}
\usepackage{geometry}
\usepackage{array}
\usepackage{tabularx}

\usepackage{hyperref}

\geometry{margin=1in}

\theoremstyle{definition}
\newtheorem{definition}{Definition}[section]
\newtheorem{theorem}{Theorem}[section]
\newtheorem{lemma}{Lemma}[section]
\newtheorem{corollary}{Corollary}[section]
\newtheorem{remark}{Remark}[section]
\newtheorem{principle}{Principle}[section]

\title{Inductive Monad over Two-Truth Adjunction: \\
From Logical Systems to Programming Languages}

\author{Masaomi Chiba}
\date{2026-02-25}

\begin{document}

\maketitle

\begin{abstract}
We introduce the \emph{inductive monad} $I_N$ as a universal mechanism that transforms any logical or grammatical system $N$ into an executable programming language. Here ``weak'' means that, on its own, a logical system does not yet carry the execution structure required of a programming language (control flow, state threading, search, and so on); this is a design gap rather than a defect, and monadic structure naturally fills it. Using the two-truth adjunction $F_n \dashv G_{2^n+1}$ (between the quotidian ($C/Dfix0$) and ultimate ($C^D$) perspectives), we establish that adding the appropriate inductive monad to any system $N$ yields a Turing-complete (or computationally meaningful) language. We present four canonical examples spanning logic, computation, and grammar—Peano arithmetic, minimal logic, context-free grammar, and Horn clauses—and provide a complete correspondence table mapping ten fundamental systems to their optimal monad instantiations. The framework unifies program design principles and provides formal evidence that ``gaps in expressiveness'' are not random but structurally determined by the missing monad.
\end{abstract}

\section{Introduction}

\subsection{The Fundamental Gap}

Every formal system exhibits a characteristic gap between \emph{description} and \emph{execution}:

\begin{itemize}
  \item Peano arithmetic can describe addition but not compute it freely.
  \item Minimal logic can express implication but not negation or control flow.
  \item Context-free grammars can specify syntax but not manage parsing state.
  \item Horn clauses describe relations but not backtracking search explicitly.
\end{itemize}

The standard response is: ``add more axioms,'' ``extend the logic,'' or ``use a richer formalism.'' But this misses a deeper pattern. We ask instead:

\begin{quote}
\emph{What if these gaps are not defects, but ``canonical missing structure''—and that structure is always a monad?}
\end{quote}

\subsection{The Main Claim}

\begin{theorem}[Informal]
For any formal system $N$ (logical, grammatical, or computational), there exists a canonical \emph{inductive monad} $I_N$ such that the composite
\[
  \text{Lang}(N) := N \otimes I_N
\]
is a programming language with well-defined execution semantics, control flow, and (under mild conditions) termination guarantees via the two-truth adjunction.
\end{theorem}

This paper makes that claim precise and demonstrates it across ten different systems.

\subsection{Related Work and Positioning}

Monads in programming are well-studied (Moggi, Wadler, Spivey). The two-truth adjunction is novel (Chiba, 2025). Our contribution is the \emph{systematic correspondence} between logical gaps and monad instantiation, unified under the two-truth framework.

---

\section{Two-Truth Adjunction and Inductive Monad}

\subsection{The Two-Truth Adjunction Recap}

\begin{definition}[Two-Truth Adjunction]
The two-truth adjunction $F_n \dashv G_{2^n+1}$ is a pair of 2-functors:
\begin{itemize}
  \item $F_n : C/D_{\mathrm{fix}0} \to C^D$ (compression: meta-level ascent, left-lax)
  \item $G_{2^n+1} : C^D \to C/D_{\mathrm{fix}0}$ (stalemate: constraint via odd cycles, right-strict)
\end{itemize}
where $C/D_{\mathrm{fix}0}$ is the quotidian (descriptive) 2-slice category and $C^D$ is the ultimate (computational) EM 2-category with Debt monad $D$.

The adjunction guarantees:
\[
  \text{Hom}_{C^D}(F_n(X), Y) \simeq \text{Hom}_{C/D_{\text{fix}0}}(X, G_{2^n+1}(Y))
\]
\end{definition}

The quotidian $C/D_{\mathrm{fix}0}$ is the realm of \emph{description}—what we wish to say. The ultimate $C^D$ is the realm of \emph{execution}—how it is computed. The adjunction ensures they remain in harmony: any gap in descriptive power is compensated by computational structure.

\subsection{Inductive Monad Definition}

\begin{definition}[Inductive Monad]
Given a logical or grammatical system $N$ (viewed as a finitary syntax or theory), an \emph{$N$-inductive monad} $I_N$ is a monad on $\mathbf{Set}$ (or a suitable category) defined such that:

\begin{enumerate}
  \item The \emph{object map} embeds the productions or rules of $N$ as generators of a free structure.
  \item The \emph{bind operation} $\mu : I_N \circ I_N \Rightarrow I_N$ realizes the ``missing control flow'' of $N$.
  \item The \emph{return} operation $\eta : \text{Id} \Rightarrow I_N$ embeds raw data as trivial programs.
\end{enumerate}

The \emph{Kleisli category} $\mathbb{K}(I_N)$ models the executable language: objects are types, arrows are programs.
\end{definition}

\begin{remark}
The inductive monad is not arbitrary. Given $N$, the monad $I_N$ is \emph{determined} by which computational effects are missing in $N$. This determination is the content of our correspondence theorems below.
\end{remark}

\subsection{Connection to Two-Truth Adjunction}

The role of the adjunction is to ensure that the iterative application of $F_n$ (meta-level ascent) and $G_{2^n+1}$ (constraint by stalemate) drives the monad's accumulated Debt (unresolved complexity) to zero in finite time.

\begin{lemma}[Monad Convergence via Adjunction]
If the inductive monad $I_N$ is paired with the quotidian/ultimate categories via the two-truth adjunction, then repeated application
\[
  F_n \circ G_{2^n+1} \circ F_{n+1} \circ G_{2^{n+1}+1} \circ \cdots
\]
causes the norm $\|\text{Debt}\|$ to decrease strictly, reaching zero in finite steps.

\emph{Proof sketch:} Each composition $F_n \circ G_{2^n+1}$ compresses syntactic depth (K-complexity) while enforcing constraint structure (\(\ell\)-completeness). The norm $\|\text{Debt}\| = K + \ell$ decreases by well-founded descent.
\end{lemma}

---

\section{Main Theorems}

\begin{theorem}[Universal Monad Correspondence]
\label{thm:universal}
For every formal system $N$ (logic, grammar, or computation), there exists a unique (up to isomorphism) inductive monad $I_N^*$ such that $N \otimes I_N^*$ exhibits the minimal computational power needed to make $N$ executable. Moreover, this monad is recoverable from the two-truth adjunction:
\[
  I_N^* = \text{colim}_{k} G_{2^k+1}(F_k(N))
\]
\end{theorem}

\begin{remark}
This theorem formalizes the intuition: the missing structure is not invented ad hoc; it emerges from the quotidian/ultimate duality itself.
\end{remark}

\begin{theorem}[Monad-Language Correspondence]
\label{thm:correspondence}
Let $N$ be a formal system and $I_N$ an inductive monad. The following are equivalent:

\begin{enumerate}
  \item The language $\text{Lang}(N, I_N) := N \otimes I_N$ is Turing-complete (or, for subrecursive $N$, achieves the intended computational class).
  \item The monad $I_N$ captures all ``missing effects'' in $N$: every step of program execution corresponds to a monad bind $\mu$.
  \item The norm $\|\text{Debt}_N(I_N)\|$ (accumulated unresolved complexity) reaches zero in finite iterations of $F_n \circ G_{2^n+1}$.
\end{enumerate}
\end{theorem}

\begin{remark}
This theorem explains why certain monad-system pairs work (Peano + Free, Horn + List) while others fail (binary logic + any monad).
\end{remark}

\begin{theorem}[Termination via Two-Truth Adjunction]
\label{thm:termination}
If a program in $\text{Lang}(N, I_N)$ is well-typed (i.e., inhabits a type in the quotidian category $C/D_{\mathrm{fix}0}$ and respects the Debt algebra of $C^D$), then its execution terminates and reaches the fixed point $D_{\mathrm{fix}0}$ (the ultimate description in the context of emptiness).

\emph{Proof sketch:} Every well-typed program corresponds to a quotidian morphism with finite norm. The adjunction $F_n \dashv G_{2^n+1}$ ensures this norm strictly decreases with each reduction step, guaranteeing termination by well-founded descent on $\mathbb{N}$.
\end{theorem}

---

\section{Four Canonical Examples}

\subsection{Example 1: Peano Arithmetic + Free Monad}

\textbf{System:} Peano arithmetic $\mathbb{P}$ with rules: $0 \in \mathbb{N}$ and $S : \mathbb{N} \to \mathbb{N}$.

\textbf{Gap:} $\mathbb{P}$ can \emph{describe} that ``$2+3=5$'' but cannot \emph{compute} it freely.

\textbf{Missing Monad:} $I_{\mathbb{P}} = \text{Free}(\Sigma_+)$ where $\Sigma_+$ is the functor with one generator per primitive operation.

\begin{definition}[Signature]
% 修正箇所: $$...$$ を \[...\] に変更し、より安全なディスプレイ数式環境にしました
\[\Sigma_+(X) = \mathbb{N} + (X \times X)\]
(either a natural number constant, or a pair of sub-expressions).
\end{definition}

\textbf{Result:} $\mathbb{P} \otimes I_{\mathbb{P}}$ becomes a \emph{term rewriting system} (or pure functional language) where:
\begin{itemize}
  \item $\text{return}(n)$ lifts a numeral into the monad.
  \item $\mu$ (bind) allows syntax trees to be composed (e.g., $(+) \otimes (\times)$ builds nested expressions).
  \item Evaluation (interpreter) unfolds the free structure via recursion schemes.
\end{itemize}

\textbf{Key insight:} The free monad's recursive structure exactly captures what Peano arithmetic lacked: the ability to \emph{suspend and compose} computations.

\subsection{Example 2: Minimal Logic + Continuation Monad}

\textbf{System:} Minimal logic $\text{ML}$ with only $(\Rightarrow, \bot)$; no explicit negation.

\textbf{Gap:} ML cannot express $\neg p$ or any form of escape/exception.

\textbf{Missing Monad:} $I_{\text{ML}} = \text{Cont}_R$ (continuation monad with answer type $R$).

\begin{definition}[Continuation Monad]
\[\text{Cont}_R(a) = (a \to R) \to R\]
\[\mu(m, k) = \lambda f. m(\lambda x. k(x)(f))\]
\[\eta(x) = \lambda f. f(x)\]
\end{definition}

\textbf{Result:} $\text{ML} \otimes I_{\text{ML}}$ becomes \emph{classical logic} with call/cc (non-local exits). Specifically:
\begin{itemize}
  \item Implication $p \Rightarrow q$ is now a \emph{program}: apply to evidence of $p$, get evidence of $q$.
  \item Contradiction $\bot$ triggers an escape (via call/cc).
  \item The continuation monad makes negation ``emerge'' from control flow.
\end{itemize}

\textbf{Key insight:} Classical logic is minimal logic plus the ability to \emph{suspend computation} (continuations). The monad reveals the hidden structure.

\subsection{Example 3: Context-Free Grammar + State Monad}

\textbf{System:} CFG rules like $S \to \text{NP} \, \text{VP}$, which are purely generative.

\textbf{Gap:} CFG cannot track \emph{position} during parsing.

\textbf{Missing Monad:} $I_{\text{CFG}} = \text{State}_{\mathbb{N}}$ where $\mathbb{N}$ is the input position.

\begin{definition}[State Monad]
\[\text{State}_S(a) = S \to (a \times S)\]
\[\mu(m, k) = \lambda s. \text{let } (x, s') = m(s) \text{ in } k(x)(s')\]
\[\eta(x) = \lambda s. (x, s)\]
\end{definition}

\textbf{Result:} $\text{CFG} \otimes I_{\text{CFG}}$ becomes a \emph{recursive descent parser} with automatic position tracking:
\begin{itemize}
  \item Each grammar rule becomes a parser function.
  \item The state carries the current position in the input.
  \item Composition of rules (via bind) automatically threads position through.
\end{itemize}

\textbf{Key insight:} Parsing requires ``memory'' (position). The state monad lifts this into the type system.

\subsection{Example 4: Horn Clauses + List Monad}

\textbf{System:} Horn clauses: $p \leftarrow q_1, \ldots, q_n$. Purely relational.

\textbf{Gap:} Horn clauses have no notion of \emph{choice points} or \emph{backtracking}.

\textbf{Missing Monad:} $I_{\text{Horn}} = \text{List}$ (with nondeterminism).

\begin{definition}[List Monad]
\[\text{List}(a) = [a]\]
\[\mu(m) = \text{concat}(m)\]
\[\eta(x) = [x]\]
\end{definition}

\textbf{Result:} $\text{Horn} \otimes I_{\text{Horn}}$ becomes \emph{Prolog}: unification + backtracking search.
\begin{itemize}
  \item Each Horn clause is a choice point (a possible solution).
  \item bind (via concatenation) explores all possibilities.
  \item Prolog's search tree is exactly the tree of list monadic choices.
\end{itemize}

\textbf{Key insight:} Logic programming is relational specification plus explicit search (via nondeterminism monad).

---

\section{Monad-System Correspondence Table}

\begin{table}[h!]
\centering
\small
\begin{tabularx}{\linewidth}{l >{\raggedright\arraybackslash}X l >{\raggedright\arraybackslash}X}
\toprule
\textbf{System } $N$ & \textbf{Gap / Missing} & \textbf{Inductive Monad } $I_N$ & \textbf{Result Language} \\
\midrule
% 修正箇所: テキストモード内にあった \Sigma や _ などを $ $ で囲みました
1. Peano Arithmetic & Composition, evaluation & $\texttt{Free}(\Sigma)$ & Arithmetic language \\
2. Simply-typed $\lambda$ & Variable binding, context & \texttt{Reader} + \texttt{Free} & Meta-programming lang. \\
3. Minimal Logic & Negation, control flow & $\texttt{Cont}_R$ & Classical logic + call/cc \\
4. Horn Clauses & Backtracking, choice & \texttt{List} / Nondet & Prolog \\
5. Regular Grammar & Enumeration, parsing & $\texttt{Free}(\Sigma)$ & Regex engine \\
6. Context-Free Grammar & Position tracking & \texttt{State} + \texttt{Free} & Recursive descent parser \\
7. Dependent Type Theory & Proof construction & $\texttt{Free}(\text{proof trees})$ & Proof assistant / script lang. \\
8. Presburger Arithmetic & Solution space search & \texttt{Writer} + \texttt{List} & Constraint solver \\
9. Untyped $\lambda$-calculus & Error recovery & \texttt{Maybe} / \texttt{Either} & Safe interpreter \\
10. Quantifier-free FOL & Nondeterministic search & \texttt{List} / Nondet & Search programming \\
\bottomrule
\end{tabularx}
\caption{Ten canonical systems and their corresponding inductive monads.}
\label{tab:correspondence}
\end{table}

\begin{remark}
Each row represents a discovery: a formal system, plus the unique monad that ``completes'' it into an executable language. The pattern is not coincidental—it is determined by the two-truth adjunction.
\end{remark}

---

\section{Meta-Level: Why These Monads?}

\subsection{The Principle of Minimal Addition}

\begin{principle}
The inductive monad $I_N$ for a system $N$ is the \emph{minimal} addition that makes the following commute:
\[
\begin{tikzcd}
N \arrow[r, "I_N"] \arrow[d] & N \otimes I_N \arrow[d, "\text{exec}"] \\
\text{Syntax} \arrow[r] & \text{Computation}
\end{tikzcd}
\]
\end{principle}

This is formalized via the adjunction: $I_N$ is the image of $N$ under $F_n$, the compression that internalizes missing structure.

\section{Meta-Level Completion: From Binary Logic to Languages}

\begin{theorem}[Meta-Completion of Binary Logic]
While pure binary logic $B$ cannot be completed by a single monad, the meta-level ascent $F_n$ automatically generates a structural gap that induces a canonical monad $I_{B^{(n)}}$.

Thus, the composite language
\[
  \text{Lang}(B^{(n)}) := B^{(n)} \otimes I_{B^{(n)}}
\]
is Turing-complete, where $B^{(n)} = F_n(B)$ and $I_{B^{(n)}}$ is the inductive monad arising from the two-truth adjunction.

In other words, the incompleteness of binary logic is not a defect but a pointer to the necessity of meta-level structure---which is precisely the space of emptiness (\emph{Śūnyatā}).
\end{theorem}

---

\section{Application: LLM Inference as Inductive Monad}

Consider an LLM's inference loop:

\begin{enumerate}
  \item Generate a token (return).
  \item Check consistency with context (bind).
  \item Either continue or stop (monad choice / effect).
\end{enumerate}

Each step is a monadic operation. The stopping condition (termination) is exactly the condition in Theorem \ref{thm:termination}: the accumulated Debt (unresolved contradictions in the reasoning chain) reaches zero.

\begin{corollary}[LLM Termination]
An LLM's inference process, when viewed as a program in $\text{Lang}(N, I_N)$ for appropriate $N$ and $I_N$, terminates if and only if its reasoning chain reaches a state with zero norm (i.e., fully grounded in the ultimate category $C^D$).
\end{corollary}

This formalizes the intuition that an LLM ``knows to stop thinking'' when it has resolved all contradictions.

---

\section{Open Problems}

\begin{enumerate}
  \item \textbf{Automated Monad Discovery:} Given a system $N$, can we algorithmically compute $I_N$? (Likely yes, via homology or Tor computations.)
  
  \item \textbf{Multi-Monad Systems:} What if a system requires multiple monads (e.g., State + Exception)? Is there a normal form?
  
  \item \textbf{Reverse Engineering:} Given a programming language $L$, can we factor it as $N \otimes I_N$ and recover the minimal system $N$?
  
  \item \textbf{Computational Complexity:} Is there a relationship between the monad's ``rank'' and the Turing degree of the resulting language?
  
  \item \textbf{Proof Support:} How does Theorem \ref{thm:termination} extend to dependent types and homotopy type theory?
\end{enumerate}

---

\section{Conclusion}

We have shown that the gap between a formal system and an executable programming language is not a defect to be patched arbitrarily, but a \emph{canonical, mathematically determined structure}—the inductive monad. The two-truth adjunction $F_n \dashv G_{2^n+1}$ explains why: the quotidian perspective (description) and the ultimate perspective (execution) are not separate; they are dual faces of a single unified computational process.

The inductive monad principle unifies seemingly disparate programming constructs (recursion, control flow, backtracking, state) under a single mathematical framework. This suggests a deep principle: \emph{expressiveness is not an accident, but an emergent property of filling the canonical gaps in a formal system}.

---

\begin{thebibliography}{99}

\bibitem{Moggi1991} E. Moggi, ``Notions of computation and monads,'' \emph{Information and Computation}, 93(1):55--92, 1991.

\bibitem{Wadler1992} P. Wadler, ``The essence of functional programming,'' \emph{Proc. POPL}, pp.~1--14, 1992.

\bibitem{MacLane1971} S. Mac Lane, \emph{Categories for the Working Mathematician}, Springer, 1971.

\bibitem{Lawvere1969} F. W. Lawvere, ``Adjointness in foundations,'' \emph{Dialectica}, 23(3-4):281--296, 1969.

\bibitem{Chiba2025} M. Chiba, ``Two-truth adjunction and Alethetic Complexity Integration Formalism,'' 2025 (forthcoming).

\bibitem{Spivey1990} M. Spivey, ``A categorical approach to the semantics of name binding,'' \emph{Theoretical Computer Science}, 151(2):437--475, 1990.

\bibitem{Plotkin2004} G. D. Plotkin and J. Power, ``Computational effects and operations,'' \emph{Proc. APPSEM II}, LNCS 3622, Springer, 2005.

\end{thebibliography}

\end{document}
